\chapter{Real numbers and sequences}
\label{sec:real_numbers_and_sequences}

\section{Informal invitation to sequences}
\label{sec:reals-intro}

I'm sure that you are probably already familiar with the concept of \textbf{sequence}. $1, 2, 3, 4, 5, \dots$ is a sequence of numbers which are increasing by one, $2, 4, 8, 16, 32, \dots$ is a sequence of numbers which are doubling, etc. If the sequence is increasing by $d$ each time, it's an arithmetic sequence. If it's multiplied by $r$ each time, it's a geometric sequence. It's a basic concept that most have already encountered in high school algebra. So, why do we have to analyze it here again?

The answer is two-fold. Firstly, we use sequences to actually construct the reals. Secondly, sequences are used to define limit of functions, which are the building blocks of derivatives. We shall answer the problem of constructing the reals in this chapter, and the limit of functions in later chapters.

The real numbers are constructed differently than all other numbers that we've constructed. We can't just introduce another operation and some equivalence class. You might say that we can just construct rational exponentiation, and we can get the reals. Yes, there are infinitely many real numbers that's obtainable by construction of new operations, such as square root or logarithm; however, there are many real numbers which cannot be constructed by algebraic operations. Those numbers are called \textbf{non-algebraic numbers} or \textbf{transcendental numbers}. Even more so, there are infinitely many of them. So our construction of the reals must be able to capture all these cases. How? The answer is through the analysis of sequences.

Firstly, let's define what a sequence is
\begin{df}{Sequence}{reals-sequence}
	A \textbf{sequence} $a_n$ of a set $\mathbb{F}$ is a mapping $a$ from $\mathbb{N}$ to $\mathbb{F}$.
\end{df}

\begin{notation}{}{reals-sequence}
	A sequence of a set $\mathbb{F}$ is written as $\sequence{a_n}{n = i}{f}$ where the starting point $i \in \mathbb{N}$ and the ending point is $f \in \mathbb{N}$. If $f = \infty$, then $\sequence{a_n}{n = i}{\infty}$ is an infinite sequence which starts at $i$.
\end{notation}

For example, a finite sequence of rationals, $\sequence{a_n}{n = 2}{7}$ defined as $a_n = \frac{1}{n}$ forms the set
\begin{equation*}
	\ab\{\frac{1}{2}, \frac{1}{3}, \frac{1}{4}, \frac{1}{5}, \frac{1}{6}, \frac{1}{7}\}.
\end{equation*}
A finite sequence of rationals, $\sequence{b_n}{n = 5}{7}$, defined as $b_n = n^2$ forms the set
\begin{equation*}
	\ab\{25, 36, 49, 64, 81, 100\}.
\end{equation*}
An infinite sequence of rationals $\sequence{c_n}{n = 3}{\infty}$ defined as $c_n = n + 2$ starting at $n = 3$ is the set
\begin{equation*}
	\ab\{5, 6, 7, 8, 9, 10, \dots\}.
\end{equation*}

Cauchy constructed the real numbers by saying that the real numbers are the \emph{limit} of these \emph{converging rational sequences}. To illustrate what \emph{limits} and \emph{converging sequences} are, let's look at some examples. Consider the sequence
\begin{gather*}
	3, 3.1, 3.14, 3.141, 3.1415, 3.14159, \dots.
\end{gather*}
This sequence, although filled with rational numbers, are \emph{gradually approaching} the real number $\cpi$, which is not rational. Consider another sequence
\begin{equation}
	1, 1.4, 1.41, 1.4142, 1.41421, \dots.
\end{equation}
This sequence is a sequence of rationals that \emph{gradually approaches} $\sqrt{2}$. These numbers here, $\cpi$ and $\sqrt{2}$, though not in the set of rationals, are the number that these sequence \emph{gradually approaches}. In math lingo, if a sequence \emph{gradually approaches} some number $L$, then that sequence is a \textbf{convergent sequence} and $L$ is said to be the \textbf{limit} of that sequence, i.e., the sequence converges to $L$. Here, you can clearly see how the reals can be possibly defined to be the limit of said converging sequence. However, we quickly run into a problem.

As you'll see later, the problem with convergence is that one has to know the limit beforehand before showing that it actually \emph{converges}. E.g., we have to know what the exact value of $\cpi$ in order to show that the sequence
\begin{equation*}
	3, 3.1, 3.14, 3.141, 3.1415, 3.14159, \dots
\end{equation*}
is actually a convergence sequence in the first place. I.e., we want to construct the reals by using the limit of a rational sequence, but we need to know what the reals are in order to show that the sequence actually converges. This is circular reasoning, and is undesirable in mathematics. To break out of this reasoning loop, we must show that a sequence of rational converges without knowing its limit in the first place.

Cauchy's constructed the reals by showing that all convergent sequences are Cauchy sequence. Informally,
\begin{displayquote}
	\emph{
		A Cauchy sequence is a sequence whose elements become arbitrarily close to each other as the sequence progresses.
	}
\end{displayquote}
After which, the algebraic operations on the reals will then be defined as the operation on each term of the sequence. E.g., if the real $x$ is defined as the limit of the sequence $\sequence{a_n}{n = i}{\infty}$ and the real $y$ as the limit of the sequence $\sequence{b_n}{n = i}{\infty}$. Their sum can then be represented as the limit of $\sequence{a_n + b_n}{n = i}{\infty}$, their product as the limit of $\sequence{a_n b_n}{n = i}{\infty}$, and their ratio as the limit of $\sequence{\frac{a_n}{b_n}}{n = i}{\infty}$, etc.

From now on, I'd like to focus on sequence of rationals because we want to construct the reals from the rationals. So, every sequence from now on is assumed to be a rational sequence. When the reals ($\mathbb{R}$) are finally constructed, one can replace all $\mathbb{Q}$ with $\mathbb{R}$.

\section{Cauchy sequences}

\begin{df}{Cauchy sequences}{reals-cauchy}
	If a rational sequence $\sequence{a_n}{n = m}{\infty}$ is Cauchy, then
	\begin{align}
		\forall (\epsilon > 0 \in \mathbb{Q}) \exists (N > m \in \mathbb{N}) \forall (j, k \geq N \in \mathbb{N}) : |a_j - a_k| \leq \epsilon.
	\end{align}
\end{df}

I understand that this definition looks very unfriendly, so let's break it down and see how each element come together to capture concept of a sequence whose elements get closer together as the sequence progresses. In my humble opinion, the approach taken by Terence Tao in his \citetitle{tao-1-2022} book \cite{tao-1-2022} is one of the most intuitive; therefore, I shall follow his teachings and deconstruct the definition of Cauchy sequences into two bite-sized ones: $\epsilon$-steadiness and eventual $\epsilon$-steadiness.

The first definition defines what it means for \emph{two elements to be close to each other.}
\begin{df}{$\epsilon$-closeness}{reals-epsilon_closeness}
	Two numbers $a$ and $b$ are $\epsilon$-close if $|a - b| \leq \epsilon$
\end{df}
This definition measures the distance on the number line. If the distance between $a$ and $b$ is less than $\epsilon$, we say that $a$ and $b$ are $\epsilon$-close.

\begin{df}{$\epsilon$-steadiness}{reals-epsilon_steadiness}
	A rational sequence $\sequence{a_n}{n = i}{f}$ is $\epsilon$-steady if for all $j, k$ between $i$ and $f$, $a_j$ and $a_k$ are $\epsilon$-close.
\end{df}

One interpretation of a sequence that's $\epsilon$-steady is that the range of the sequence ($\max\ab(\sequence{a_n}{n = i}{f}) - \min\ab(\sequence{a_n}{n = i}{f})$) is restricted to be less than or equal to $\epsilon$:
\begin{equation}
	\max\ab(\sequence{a_n}{n = i}{f}) - \min\ab(\sequence{a_n}{n = i}{f}) \leq \epsilon.
\end{equation}
Because of this, the distance between each term in an $\epsilon$-steady sequence is always less than or equal to $\epsilon$.

This definition also works for infinite sequence. However, the concept of maxima and minima is currently a bit loose. Sometimes, the maxima and minima doesn't even exist, so it'd be wiser for now to think of an $\epsilon$-steady sequence as a sequence that the distance between each and every term is less than $\epsilon$. Now, let's look at some examples and non-examples.

\begin{exmp}{$\epsilon$-steady sequences}{}
	\begin{enumerate}
		\item The sequence $\sequence{a_n}{n = 0}{\infty}$: $3, 3.1, 3.14, 3.141, 3.1415, 3.14159, \dots$, is $0.1$-steady but not $0.01$-steady, because
		      \begin{equation*}
			      |a_0 - a_1| = |3 - 3.1| = 0.1, \quad \textrm{which is more than} \quad 0.01.
		      \end{equation*}
		\item The sequence $\sequence{b_n}{n = 0}{\infty}$: $-1, 0, 1, 0, -1, 0, \dots$ is $2$-steady, but not $1$-steady because
		      \begin{equation*}
			      |b_0 - b_2| = |-1 - 1| = 2, \quad \textrm{which is more than} \quad 1.
		      \end{equation*}
		\item The sequence $\sequence{c_n}{n = 4}{9}$: $2, 3, 1, 5, 4, 6$ is $10$-steady, but not $4$-steady because
		      \begin{equation*}
			      |c_6 - c_9| = |6 - 1| = 5, \quad \textrm{which is more than} \quad 4.
		      \end{equation*}
	\end{enumerate}
\end{exmp}

Before moving on, I'd like to discuss the image of $\epsilon$-steady sequence. Consider a sequence $\sequence{a_n}{n = 0}{\infty}$ shown in \cref{fig:reals-epsilon_steady_sequences}. Let's draw a rectangle which infinite extends along the $x$-axis but has the height of $\epsilon$. The sequence is $\epsilon$-steady \ifft you can find a position of this rectangle such that every term in the sequence lies in that rectangle. Here, $\sequence{a_n}{n = 0}{\infty}$ is $\epsilon$-steady because there's a rectangle configuration in which every point of the sequence lies in the rectangle shown in the right most figure. The left most and middle figure doesn't count because there are some points, drawn in red, which is outside the rectangle.

\begin{figure}[]
	\centering
	\includegraphics{realanalysis/realsandsequences/cauchysequences/epsilon-steady-sequences.pdf}
	\caption{Examples of rectangle configurations to check that a sequence is $\epsilon$-steady.}
	\label{fig:reals-epsilon_steady_sequences}
\end{figure}

The concept of a sequence being $\epsilon$-steady is quite simple, but it doesn't capture the limiting behavior of a sequence. Since we want to construct the reals from the \emph{limit} of a \emph{converging sequence}, we shouldn't really care about the initial terms. Therefore, we improve on the definition of $\epsilon$-steadiness to eventual $\epsilon$-steadiness: a definition that captures what it means for \emph{the elements of a sequence to get close together as the sequence progresses}.
\begin{df}{Eventual $\epsilon$-steadiness}{reals-eventual_epsilon_steadiness}
	A rational sequence $\sequence{a_n}{n = i}{f}$ is eventually $\epsilon$-steady if there exists some natural $N \geq i$ in which the sequence $\sequence{a_n}{n = N}{f}$ is $\epsilon$-steady.
\end{df}

The intuition for this is that, if it's possible to find some number $N$ in which the removal of the first $N$ terms of the sequence makes the sequence $\epsilon$-steady, then the sequence is eventually $\epsilon$-steady. Let's look at some examples and non-examples.
\begin{exmp}{Eventually $\epsilon$-steady sequences}{reals-eventual_epsilon_steadiness}
	\begin{enumerate}
		\item The sequence $\sequence{a_n}{n = 1}{\infty}$: $1, \frac{1}{2}, \frac{1}{3}, \frac{1}{4}, \dots$, is not $\frac{1}{2}$-steady because
		      \begin{equation*}
			      |a_0 - a_2| = \ab|1 - \frac{1}{3}| = \frac{2}{3}, \quad \textrm{which is more than} \quad \frac{1}{2};
		      \end{equation*}
		      however, the sequence is eventually $\frac{1}{2}$-steady because the removal of the first term yields a sequence that is $\frac{1}{2}$-steady:
		      \begin{equation*}
			      \frac{1}{2}, \frac{1}{3}, \frac{1}{4}, \frac{1}{5}, \dots,
		      \end{equation*}
		\item The sequence $\sequence{b_n}{n = 0}{\infty}$: $1, \frac{1}{2}, \frac{1}{4}, \frac{1}{8}, \dots$, is not $\frac{1}{4}$-steady because
		      \begin{equation*}
			      |a_2 - a_4| = \ab|\frac{1}{2} - \frac{1}{8}| = \frac{3}{8}, \quad \textrm{which is more than} \quad \frac{1}{4};
		      \end{equation*}
		      however, the sequence is eventually $\frac{1}{4}$-steady because the removal of the first two terms yields a sequence that is $\frac{1}{4}$-steady:
		      \begin{equation*}
			      \frac{1}{4}, \frac{1}{8}, \frac{1}{16}, \frac{1}{32}, \dots.
		      \end{equation*}
		\item The sequence $\sequence{c_n}{n = 0}{\infty}$: $69, 2, 2, 2, 2, \dots$ is not $3$-steady because for all $j > 0$,
		      \begin{equation*}
			      |a_0 - a_j| = |69 - 2| = 67, \quad \textrm{which is more than} \quad 3;
		      \end{equation*}
		      however, the sequence is eventually $3$-steady because the removal of the first term yields a sequence that is $3$-steady:
		      \begin{equation*}
			      2, 2, 2, 2, 2, \dots.
		      \end{equation*}
	\end{enumerate}
\end{exmp}

And now, we're ready to define what Cauchy sequences are in simple terms:
\begin{cor}{Cauchy sequences}{reals-cauchy_sequences_redefined}
	A Cauchy sequence is an infinite sequence which is eventually $\epsilon$-steady for all $\epsilon > 0 \in \mathbb{Q}$. I.e., for a sequence $\sequence{a_n}{n = m}{\infty}$ to be Cauchy,
	\begin{equation*}
		\forall (\epsilon > 0 \in \mathbb{Q}) \exists (N > m \in \mathbb{N}) \forall (j, k \geq N \in \mathbb{N}) : |a_j - a_k| \leq \epsilon
	\end{equation*}
\end{cor}

A natural question to ask from here would be, \emph{how to prove that a sequence is Cauchy?} It is quite hard, but the core concept is to show that for every $\epsilon$, there exists a corresponding $N$ which makes the sequence $\sequence{a_n}{n = N}{\infty}$, $\epsilon$-steady. This would be one of the first proof that must have a preparation stage before the actual proof. We do some arithmetic first to decide what $N$ might be in terms of $\epsilon$, and use that to show that for all $j, k \geq N$, the sequence is $\epsilon$-steady.

\begin{exmp}{Show that $\sequence{a_n}{n = 0}{\infty}$ is Cauchy}{}
	We prepare the proof by first fixing $\epsilon$ and find $N$ in terms of $\epsilon$. For all $n \geq N$, we demand that
	\begin{align}
		|a_j - a_k|                         & \leq \epsilon                                    \\
		\ab|\frac{1}{j} - \frac{1}{k}|      & \leq \epsilon                                    \\
		\ab|\frac{1}{j}| + \ab|\frac{1}{k}| & \leq \epsilon. &  & \textrm{Triangle inequality}
	\end{align}
	If we demand both $\ab|\frac{1}{j}|$ and $\ab|\frac{1}{k}|$ to be less than $\frac{\epsilon}{2}$, then we would get $|a_j - a_k| \leq \epsilon$ as desired; thus, let
	\begin{equation}
		\ab|\frac{1}{j}|, \ab|\frac{1}{k}| \leq \frac{\epsilon}{2}
	\end{equation}
	Since $j, k \geq N$, we have $\frac{1}{j}, \frac{1}{k} \leq \frac{1}{N}$. Combining this with the equation above gives
	\begin{align}
		\frac{1}{N} & \leq \frac{\epsilon}{2}  \\
		N           & \geq \frac{2}{\epsilon}.
	\end{align}
	Here, we've shown that $N \geq \frac{2}{\epsilon}$ causes $|a_j - a_k|$ where $j, k \geq N$ to be less than $\epsilon$ for every $\epsilon$; thus, the sequence is Cauchy. Now onto the formal proof.

	\begin{proof}
		\begin{align}
			|a_j - a_k|                             & = \ab|\frac{1}{j} - \frac{1}{k}|                                                 \\
			\ab|\frac{1}{j} - \frac{1}{k}|          & \leq \ab|\frac{1}{j}| + \ab|\frac{1}{k}|     &  & \textrm{Triangle's inequality} \\
			\ab|\frac{1}{j}| + \ab|\frac{1}{k}|     & \leq \frac{1}{N} + \frac{1}{N}                                                   \\
			\shortintertext{Choose $N \geq \frac{2}{\epsilon}$, then}
			\frac{1}{N} + \frac{1}{N}               & \geq \frac{\epsilon}{2} + \frac{\epsilon}{2}                                     \\
			\frac{\epsilon}{2} + \frac{\epsilon}{2} & = \epsilon.
		\end{align}
		Therefore, $N \geq \frac{2}{\epsilon}$ causes $|a_j - a_k|$ to be less than $\epsilon$ for all $j, k \geq N$. Thus, the sequence is Cauchy.
	\end{proof}
\end{exmp}

Before constructing the reals, I'd like to discuss the notion of boundedness which is later used to prove that multiplication on the reals is well-defined

\begin{df}{Boundedness of sequences}{reals-boundedness}
	A sequence $\sequence{a_n}{n = i}{f}$ is bounded by a rational $M > 0$ if for all $n \in \mathbb{N}$ that's between $i$ and $f$, $|a_n| \leq M$.
\end{df}

\begin{cor}{Finite sequences are bounded}{reals-finite_sequences_bounded}
	Every finite sequence $\sequence{a_n}{n = i}{f}$ is bounded
\end{cor}
\begin{proof}
	We shall use strong induction to prove this.

	\basecase Trivially, the sequence $a_1$ is bounded by $M = |a_1|$.

	\inductivestep Assume that the sequence $\sequence{a_n}{n = i}{N}$ where $i < N < f$ is bounded by some rational $M$. Then, the sequence $\sequence{a_n}{n = i}{S(N)}$ must also be bounded by the rational $M + |a_{S(N)}|$; therefore, any finite sequence is bounded, closing the induction.
\end{proof}

The next thing is to establish a connection between Cauchy sequences and the concept of boundedness.

\begin{prop}{Every Cauchy sequence is bounded}{reals-boundedness-cauchy}
	Every Cauchy sequence $\sequence{a_n}{n = i}{\infty}$ is bounded.
\end{prop}

\begin{proof}
	Since the sequence is Cauchy, for every $\epsilon$, there exists a natural $N$ such that for all $j, k \geq N$, $|a_j - a_k| \leq \epsilon$. The sequence can then be split into two parts:
	\begin{equation*}
		a_1, a_2, \dots, a_{N - 1}, \mathand a_{N}, a_{N + 1}, a_{N + 2}, \dots.
	\end{equation*}
	By \cref{co:reals-finite_sequences_bounded}, the first part of the sequence is bounded by some rational $M$. Since the distance between every term in the second part is less than $\epsilon$. If the smallest term in the second part of the sequence is $K$, then the entire part must be bounded by $K + \epsilon$; therefore, the whole sequence is bounded by $M + K + \epsilon$.
\end{proof}

\section{Construction of the reals}

Now we are ready to construct the reals. Similar to how we've constructed other objects, we do not care what the object actually is, but rather, how it interacts with others. Therefore, let's define a function that takes in a Cauchy sequence and spits out a real number:
\begin{df}{Real numbers}{reals-real}
	A real number is of the form $\Zreal{a_n}$ where $\sequence{a_n}{n = m}{\infty}$ is a Cauchy sequence.
\end{df}
Similar to how integers are defined, there are many Cauchy sequences which represent the same real number. I.e.,
\begin{equation}
	1, \frac{1}{2}, \frac{1}{4}, \frac{1}{8}, \frac{1}{16}, \dots, \mathand 1, \frac{1}{2}, \frac{1}{3}, \frac{1}{4}, \frac{1}{5}, \dots, \footnote{It is left as an exercise for the reader to prove that these two sequences are Cauchy.}
\end{equation}
both of which seems to approach the number $0$. Therefore, we must define the concept of equivalence sequences:
\begin{df}{Equivalent sequences}{reals-equivalent_sequences}
	Two sequences, $\sequence{a_n}{n = m_a}{\infty}$ and $\sequence{b_n}{n = m_b}{\infty}$ are equivalent \ifft for all rational $\epsilon$, there exists a natural $N \geq \max(m_a, m_b)$ in which for all $n \geq N$, $a_n$ and $b_n$ are $\epsilon$-close, i.e.,
	\begin{multline}
		\forall (\epsilon > 0 \in \mathbb{Q}) \exists \ab(N \geq \max(m_a, m_b) \in \mathbb{N}) \forall (n \geq N \in \mathbb{N}) : \\ |a_n - b_n| \leq \epsilon.
	\end{multline}
	$\sequence{a_n}{n = m_a}{\infty}$ and $\sequence{b_n}{n = m_b}{\infty}$ represent the same number, i.e., $\Zreal{a_n} = \Zreal{b_n}$ \ifft they are both Cauchy and equivalent
\end{df}

\begin{notation}{Equivalent sequences}{reals-equivalent_sequences}
	If $\sequence{a_n}{n = m_a}{\infty}$ and $\sequence{b_n}{n = m_b}{\infty}$ are equivalent, we write
	\begin{equation*}
		\sequence{a_n}{n = m_a}{\infty} \sim \sequence{b_n}{n = m_b}{\infty}.
	\end{equation*}
\end{notation}

The first thing to notice is that the notion of equivalent (\cref{df:reals-equivalent_sequences}) doesn't care about the starting point of the sequences. We analyze the equivalency for $n \geq \max(m_a, m_b)$ only, because we only care about the behavior of the sequence as it \emph{progresses}. Similar to how we've defined Cauchy sequences, we can also break \cref{df:reals-equivalent_sequences} into two parts: $\epsilon$-close, and eventual $\epsilon$-close. I shall go over them very briefly as it is analogous to how Cauchy sequences (\cref{df:reals-cauchy}) are defined.

Two sequences $\sequence{a_n}{n = m_a}{\infty}$ and $\sequence{b_n}{n = m_b}{\infty}$ are $\epsilon$-close \ifft
\begin{equation}
	|a_n - b_n| \leq \epsilon
\end{equation}
These sequences are eventually $\epsilon$-close if there exists a number $N > \max(m_a, m_b)$ in which $\sequence{a_n}{n = N}{\infty}$ and $\sequence{b_n}{n = N}{\infty}$ are $\epsilon$-close. Finally, these sequences are equivalence if it is eventually $\epsilon$-close for all rational $\epsilon$.

The second thing to notice is that \cref{df:reals-equivalent_sequences} doesn't care about whether the sequences are Cauchy or not. E.g.,
\begin{equation}
	1, 2, 3, 4, 5, 6, \dots, \mathand 3, 3, 3, 4, 5, 6, \dots,
\end{equation}
are equivalent but not Cauchy.\footnote{It is left as an exercise for the reader to prove that these two sequence aren't Cauchy, but are still equivalent.} But if two sequences are equivalent and one of them is Cauchy, we can guarantee that the other one must also be Cauchy.
\begin{cor}{}{}
	Let $\sequence{a_n}{n = m_a}{\infty} \sim \sequence{b_n}{n = m_b}{\infty}$. If one of the sequence is Cauchy, then the other must also be. If one of the sequence isn't Cauchy, then the other is also not Cauchy.
\end{cor}

\begin{proof}
	By \cref{df:reals-equivalent_sequences}, $\sequence{a_n}{n = m_a}{\infty} \sim \sequence{b_n}{n = m_b}{\infty}$ means
	\begin{equation}
		\forall (\epsilon > 0 \in \mathbb{Q}) \exists \ab(N \geq \max(m_a, m_b) \in \mathbb{N}) \forall (n \geq N \in \mathbb{N}) : |a_n - b_n| \leq \epsilon.
	\end{equation}
	Let's assume that $\sequence{a_n}{n = m_a}{\infty}$ is Cauchy, i.e.,
	\begin{equation}
		\forall (\epsilon > 0 \in \mathbb{Q}) \exists \ab(N \geq \max(m_a, m_b) \in \mathbb{N}) \forall (j, k \geq N \in \mathbb{N}) : |a_j - a_k| \leq \epsilon. \label{eq:reals-equivalence_sequences_cauchy-1}
	\end{equation}
	We then want to prove that $\sequence{b_n}{n = m_b}{\infty}$ is also Cauchy, i.e.,
	\begin{equation*}
		\forall (\delta > 0 \in \mathbb{Q}) \exists \ab(N \geq \max(m_a, m_b) \in \mathbb{N}) \forall (j, k \geq N \in \mathbb{N}) : |b_j - b_k| \leq \delta. \label{eq:reals-equivalence_sequences_cauchy-2}
	\end{equation*}
	This may seem a bit outrageous, but notice that
	\begin{equation}
		|b_j - b_k| = |b_j - a_j + a_j - a_k + a_k - b_k|.
	\end{equation}
	By the triangle's inequality, we get
	\begin{equation}
		|b_j - a_j + a_j - a_k + a_k - b_k| \leq |b_j - a_j| + |a_j - a_k| + |a_k - b_k|.
	\end{equation}
	From \cref{eq:reals-equivalence_sequences_cauchy-1,eq:reals-equivalence_sequences_cauchy-2},
	\begin{equation}
		|b_j - a_j| + |a_j - a_k| + |a_k - b_k| \leq 3 \epsilon,
	\end{equation}
	or $|b_j - b_k| \leq 3\epsilon$. If $\delta = 3\epsilon$, then $|b_j - b_k| \leq \delta$ as desired.
\end{proof}

\section{Addition and multiplication on the reals}

As said earlier, we want to define arithmetic on the reals as operations on Cauchy sequences. Since the definition of division has much nuance, we shall focus on addition, subtraction, and multiplication for now. Similar to how addition and multiplication on lower number systems, we shall define the operations, then prove that they are well-defined.

\begin{df}{Addition on the reals}{reals-addition}
	Let there be two Cauchy sequences $\sequence{a_n}{n = m_a}{\infty}$ which represents the real $x = \Zreal{a_n}$, and $\sequence{b_n}{n = m_b}{\infty}$ which represents the real $y = \Zreal{b_n}$. The sum of $x$ and $y$, denoted $x + y$, is given by the real number $\Zreal{a_n + b_n}$ represented by the sequence $\sequence{a_n + b_n}{n = \max(m_a, m_b)}{\infty}$, i.e.,
	\begin{equation}
		x + y = \Zreal{a_n} + \Zreal{b_n} = \Zreal{a_n + b_n}.
	\end{equation}
\end{df}

To actually prove that addition on the reals is well-defined, we have to make sure that the sum of two Cauchy sequences are Cauchy first.
\begin{lemma}{The sum of two Cauchy sequences are Cauchy}{reals-addition_cauchy}
	Let there be two Cauchy sequences $\sequence{a_n}{n = m_a}{\infty}$ and $\sequence{b_n}{n = m_b}{\infty}$. Then, $\sequence{a_n + b_n}{n = \max(m_a, m_b)}{\infty}$ is also a Cauchy sequence.
\end{lemma}

\begin{proof}
	$\sequence{a_n}{n = m_a}{\infty}$ and $\sequence{b_n}{n = m_b}{\infty}$ are Cauchy means that for all rational $\epsilon$,
	\begin{gather}
		\exists (A > m_a \in \mathbb{N}) \forall (j, k \geq A \in \mathbb{N}) : |a_j - a_k| \leq \epsilon, \\
		\exists (B > m_b \in \mathbb{N}) \forall (j, k \geq A \in \mathbb{N}) : |b_j - b_k| \leq \epsilon.
	\end{gather}
	Let $C = \max(A, B)$; therefore,
	\begin{equation}
		\forall (j, k \geq C \in \mathbb{N}) : \ab(|a_j - a_k| \leq \epsilon) \land \ab(|b_j - b_k| \leq \epsilon). \label{eq:reals-addition_cauchy-1}
	\end{equation}
	We want to show that $\sequence{a_n + b_n}{n = \max(m_a, m_b)}{\infty}$ is Cauchy, i.e.,
	\begin{multline}
		\forall (\delta > 0 \in \mathbb{Q}) \exists \ab(N > \max(m_a, m_b)) \forall (j, k \geq C \in \mathbb{N}) : \\ \ab|(a_j + b_j) - (a_k + b_k)| \leq \delta
	\end{multline}
	From \cref{eq:reals-addition_cauchy-1}, we get that for all $j, k > C$,
	\begin{align}
		|a_j - a_k| + |b_j - b_k|      & \leq 2\epsilon  \\
		\shortintertext{By the triangle's inequality,  $|a_j - a_k + b_j - b_k| \leq |a_j - a_k| + |b_j - b_k|$; thus,}
		\ab|a_j - a_k + b_j - b_k|     & \leq 2\epsilon  \\
		\ab|(a_j + b_j) - (a_k + b_k)| & \leq 2\epsilon.
	\end{align}
	Let $\delta = 2\epsilon$, and we get $\ab|(a_j + b_j) - (a_k + b_k)| \leq \delta$ as desired.
\end{proof}
The direct consequence of \cref{le:reals-addition_cauchy} is that the product of two Cauchy sequence is also bounded because Cauchy sequences are bounded. Then, we prove that addition on the reals is well-defined.

\begin{prop}{Addition on the reals is well-defined}{reals-addition_well_defined}
	Let there be two Cauchy sequences in which $\sequence{a_n}{n = m_a}{\infty} \sim \sequence{a_n'}{n = m_{a'}}{\infty}$, and another Cauchy sequence $\sequence{b_n}{n = m_b}{\infty}$. If $x = \Zreal{a_n}$, $x' = \Zreal{a_n'}$, and $y = \Zreal{b_n}$ then
	\begin{equation}
		x + y = x' + y
	\end{equation}
	One other way to write this proposition is that sums of equivalent Cauchy sequences are equivalent.
\end{prop}

\begin{proof}
	Let $m = \max(m_a, m_{a'}, m_b)$, then $\sequence{a_n}{n = m}{\infty} \sim \sequence{a_n'}{n = m}{\infty}$ can be written as
	\begin{equation}
		\exists (N \in \mathbb{N}) \forall (n \geq N \in \mathbb{N}) : |a_n - a_n'| \leq \epsilon. \label{eq:reals-addition_well_defined-4}
	\end{equation}
	We want to show that $x + y = x' + y$, i.e., $\sequence{a_n + b_n}{n = m}{\infty} \sim \sequence{a_n' + b_n}{n = m}{\infty}$. By \cref{le:reals-addition_cauchy}, these two sequences must be Cauchy, so the final result is guaranteed to be a real number. This proof uses the same trick as the one before: adding and subtracting the same term. From \cref{eq:reals-addition_well_defined-4},
	\begin{align}
		|a_n - a_n'|                    & \leq \epsilon  \\
		\ab|a_n - b_n + b_n - a_n'|     & \leq \epsilon  \\
		\ab|(a_n - b_n) - (a_n' + b_n)| & \leq \epsilon,
	\end{align}
	meaning that $\sequence{a_n + b_n}{n = m}{\infty} \sim \sequence{a_n' + b_n'}{n = m}{\infty}$; thus,
	\begin{align}
		\Zreal{a_n + b_n}         & = \Zreal{a_n' + b_n}.        \\
		\shortintertext{By \cref{df:reals-addition},}
		\Zreal{a_n} + \Zreal{b_n} & = \Zreal{a_n'} + \Zreal{b_n} \\
		x + y                     & = x' + y \qedhere
	\end{align}
\end{proof}

Because I do not want to write separate proof for addition and subtraction, I shall embed the concept of negative numbers in the reals and define $x - y = x + (-y)$. This way, we only have to introduce the negation of the real numbers, and prove that it is well-defined instead of defining subtraction as its own operation, which would be more cumbersome.
\begin{df}{Negation of a real number}{reals-negative}
	Given a Cauchy sequence $\sequence{a_n}{n = m}{\infty}$ in which $x = \Zreal{a_n}$, we define
	\begin{equation}
		- x = \Zreal{-a_n}.
	\end{equation}
\end{df}

\begin{prop}{Negation of a real is well-defined}{reals-negative_well_defined}
	If two Cauchy sequences $\sequence{a_n}{n = m_a}{\infty} \sim \sequence{a_n'}{n = m_{a'}}{\infty}$, and $x = \Zreal{a_n}$ and $x' = \Zreal{a_n'}$. Then,
	\begin{equation}
		- x = - x'
	\end{equation}
	Another way to say this is that the negation of two equivalent Cauchy sequences are equivalent.
\end{prop}

\begin{proof}
	$\sequence{a_n}{n = m_a}{\infty} \sim \sequence{a_n'}{n = m_{a'}}{\infty}$ means for all rational $\epsilon$,
	\begin{equation}
		\ab(N \geq \max(m_a, m_b) \in \mathbb{N}) \forall (n \geq N \in \mathbb{N}) : |a_n - a_n'| \leq \epsilon.
	\end{equation}
	From the property of absolute values, $|a_n - a_n'| = \ab|(-a_n') - (-a_n)|$; therefore, $\sequence{-a_n}{n = m_a}{\infty} \sim \sequence{-a_n'}{n = m_{a'}}{\infty}$ and $x = -x'$ as desired.
\end{proof}

Subtraction is then defined as addition of a negated real number.
\begin{df}{Subtraction on the reals}{reals-subtraction}
	If $x = \Zreal{a_n}$ and $y = \Zreal{b_n}$, then the difference between $x$ and $y$ is
	\begin{equation*}
		x - y = x + (-y) = \Zreal{a_n} + \Zreal{-b_n}.
	\end{equation*}
\end{df}
Since both negation and addition is well-defined, and negation is written using addition, subtraction must also be well-defined on the reals. We then move on to define multiplication.

\begin{df}{Multiplication on the reals}{reals-product_well_defined}
	Let there be two Cauchy sequences, $\sequence{a_n}{n = m_a}{\infty}$ and $\sequence{b_n}{n = m_b}{\infty}$, in which $x = \Zreal{a_n}$ and $y = \Zreal{b_n}$. The product of $x$ and $y$, denoted $xy$, is given by the real number $\Zreal{a_n b_n}$ represented by the sequence $\sequence{a_n b_n}{n = \max(m_a, m_b)}{\infty}$, i.e.,
	\begin{equation}
		xy = \Zreal{a_n}\Zreal{b_n} = \Zreal{a_n b_n}.
	\end{equation}
\end{df}
Similarly, we must also prove that the product of two Cauchy sequences are Cauchy:
\begin{lemma}{The product of two Cauchy sequences are Cauchy}{reals-product_cauchy}
	Let there be two Cauchy sequences $\sequence{a_n}{n = m_a}{\infty}$ and $\sequence{b_n}{n = m_b}{\infty}$. Then, $\sequence{a_n b_n}{n = \max(m_a, m_b)}{\infty}$ is also a Cauchy sequence.
\end{lemma}
\begin{proof}
	$\sequence{a_n}{n = m_a}{\infty}$ and $\sequence{b_n}{n = m_b}{\infty}$ are Cauchy means that for all rational $\epsilon$,
	\begin{gather}
		\exists (A > m_a \in \mathbb{N}) \forall (j, k \geq A \in \mathbb{N}) : |a_j - a_k| \leq \epsilon, \\
		\exists (B > m_b \in \mathbb{N}) \forall (j, k \geq A \in \mathbb{N}) : |b_j - b_k| \leq \epsilon.
	\end{gather}
	Let $C = \max(A, B)$; therefore,
	\begin{equation}
		\forall (j, k \geq C \in \mathbb{N}) : \ab(|a_j - a_k| \leq \epsilon) \land \ab(|b_j - b_k| \leq \epsilon). \label{eq:reals-multiplication_cauchy-1}
	\end{equation}
	We want to show that $\sequence{a_n b_n}{n = \max(m_a, m_b)}{\infty}$ is Cauchy, i.e.,
	\begin{equation}
		\forall (\delta > 0) \exists (N \geq C \in \mathbb{N}) \forall (j, k \geq N \in \mathbb{N}) : |a_j b_j - a_k b_k| \leq \epsilon.
	\end{equation}
	Notice that
	\begin{align}
		|a_j b_j - a_k b_k|                 & = |a_j b_j - a_k b_j + a_k b_j - a_k b_k|  \\
		                                    & = |b_j (a_j - a_k) + a_k (b_j - b_k)|      \\
		|b_j (a_j - a_k) + a_k (b_j - b_k)| & \leq |b_j| |a_j - a_k| + |a_k| |b_j - b_k| \\
	\end{align}
	From \cref{eq:reals-multiplication_cauchy-1},
	\begin{equation}
		|b_j| |a_j - a_k| + |a_k| |b_j - b_k| \leq \epsilon \ab(|b_j| + |a_k|).
	\end{equation}
	It is quite tempting to let $\delta = \epsilon \ab(|b_j| + |a_k|)$, but it'd be better if the R.H.S. is dependent on the property of the sequence as a whole and not restricted to individual terms. Here is where boundedness comes in. From \cref{pr:reals-boundedness-cauchy}, every Cauchy sequences are bounded. So, let's make $U$ and $V$ be the bound of $\sequence{a_n}{n = m_a}{\infty}$ and $\sequence{b_n}{n = m_b}{\infty}$ respectively, i.e.,
	\begin{equation}
		\forall (n \geq C \in \mathbb{N}) : \ab(|a_n| \leq U) \land \ab(|b_n| \leq V);
	\end{equation}
	thus, $\epsilon \ab(|b_j| + |a_k|) \leq \epsilon (U + V)$. If $\delta = \epsilon (U + V)$, we get $|a_j b_j - a_k b_k| \leq \delta$ as desired.
\end{proof}
The direct consequence of \cref{le:reals-product_cauchy} is, the product of two Cauchy sequence is also bounded. Then, we prove that multiplication on the reals is well-defined.

\begin{prop}{Multiplication on the reals is well-defined}{pr:reals-product_well_defined}
	Let there be two Cauchy sequences in which $\sequence{a_n}{n = m_a}{\infty} \sim \sequence{a_n'}{n = m_{a'}}{\infty}$, and another Cauchy sequence $\sequence{b_n}{n = m_b}{\infty}$. If $x = \Zreal{a_n}$, $x' = \Zreal{a_n'}$, and $y = \Zreal{b_n}$ then
	\begin{equation}
		xy = x'y
	\end{equation}
	One other way to write this proposition is that the product of equivalent Cauchy sequences are equivalent.
\end{prop}

\begin{proof}
	Let $C = \max(m_a, m_{a'}, m_b)$, then $\sequence{a_n}{n = m_a}{\infty} \sim \sequence{a_n'}{n = m_{a'}}{\infty}$ means that for all rational $\epsilon$,
	\begin{equation}
		\exists (N \geq M \in \mathbb{N}) \forall (n \geq N \in \mathbb{N}) : |a_n - a_n'| \leq \epsilon.
	\end{equation}
	Because all the sequences are Cauchy,
	\begin{multline}
		\forall (\epsilon > 0 \in \mathbb{Q}) \exists (N \geq C) \forall (j, k \geq N \in \mathbb{N}) : \\
		\ab(\ab|b_j - b_k| \leq \epsilon.) \land \ab(\ab|a_j - a_k| \leq \epsilon) \land \ab(\ab|a_j' - a_k'| \leq \epsilon).
	\end{multline}
	We want to show that $\sequence{a_n b_n}{n = C}{\infty} \sim \sequence{a_n' b_n}{n = C}{\infty}$, i.e.,
	\begin{equation}
		\forall (\delta > 0 \in \mathbb{Q}) \exists (M \geq C) \forall (m \geq M \in \mathbb{N}) : |a_m b_m - a_m' b_m| \leq \delta.
	\end{equation}
	First,
	\begin{equation}
		|a_m b_m - a_m' b_m| = |b_m||a_m - a_m'|
	\end{equation}
	\begin{itemize}
		\item From $\sequence{a_n}{n = m_a}{\infty} \sim \sequence{a_n'}{n = m_{a'}}{\infty}$, $|a_m - a_m'| \leq \delta$.
		\item Because $\sequence{b_n}{m = C}{\infty}$ is Cauchy, it is bounded by some rational $B$; thus, $\forall (m \leq C) : |b_m| \leq B$.
	\end{itemize}
	Combining these together we get
	\begin{equation}
		|b_m||a_m - a_m'| \leq B\epsilon.
	\end{equation}
	If $\delta = B\epsilon$, then $|a_m b_m - a_m' b_m| \leq \delta$ as desired.
\end{proof}

\section{Reciprocation and division on the reals}

Since division is just multiplication by the reciprocal, we shall only focus on reciprocation for now.

The reciprocal of $\sequence{a_n}{n = m}{\infty}$ is defined \ifft,
\begin{enumerate}
	\item The sequence $\sequence{a_n}{n = m}{\infty}$ is \emph{not} equivalent to the zero sequence $\sequence{0}{n = m}{\infty}$: any sequences which are equivalent to the zero sequence represents zero, and the reciprocal of zero is not defined.
	\item None of the terms in the sequence are zero. If one of the term in the sequence is zero, say $a_j = 0$, then $\frac{1}{a_j}$ is $\infty$. Since $\infty$ isn't a rational number, the sequence $\sequence{\frac{1}{a_n}}{n = m}{\infty}$ is not a rational sequence, thus by \cref{df:reals-real}, it is not a real number.
\end{enumerate}
These two restrictions forces us to define what sequence is allowed to be reciprocated. If a sequence doesn't follow these criteria, then its reciprocation is either not a Cauchy sequence of rationals; thus, can't be used to represent real numbers. Let's see some examples.
\begin{exmp}{Sequences which cannot be reciprocated}{}
	\begin{enumerate}
		\item Consider the sequence $\sequence{a_n}{n = m_a}{\infty}$
		      \begin{equation}
			      1, \frac{1}{2}, \frac{1}{3}, \frac{1}{4}, \frac{1}{5}, \dots
		      \end{equation}
		      The reciprocation of this sequence, $\sequence{\frac{1}{a_n}}{n = m_a}{\infty}$ is the sequence
		      \begin{equation}
			      1, 2, 3, 4, 5, \dots,
		      \end{equation}
		      which is obviously not Cauchy. This is because the original sequence is actually equivalent to the zero sequence $\sequence{0}{n = m}{\infty}$.
		\item Consider the sequence $\sequence{b_n}{n = m_b}{\infty}$:
		      \begin{equation}
			      0, 0.9, 0.99, 0.999, 0.9999, \dots
		      \end{equation}
		      The reciprocation of this sequence, $\sequence{\frac{1}{b_n}}{n = m_b}{\infty}$ is the sequence
		      \begin{equation}
			      \infty, 1.11111\dots, 1.01010\dots, 1.00100\dots, 1.00010\dots, \dots.
		      \end{equation}
		      Even though this looks like a Cauchy sequence that's somehow approaching the number one, the sequence itself isn't Cauchy because $\infty$ is not a number. This is because there are $0$ in the original sequence.
	\end{enumerate}
\end{exmp}
This problem leads us to define a type of sequence that can be used for reciprocation: \emph{sequences bounded away from zero}.
\begin{df}{Sequences bounded away from zero}{reals-bounded_away_from_zero}
	A sequence $\sequence{a_n}{n = m}{\infty}$ is bounded away from zero \ifft there exists some rational $c$ in which for all natural $n \geq m$, $| a_n | \geq c$, i.e.,
	\begin{equation}
		\exists (c > 0 \in \mathbb{Q}) \forall (n \geq m \in \mathbb{N}) : | a_n | \geq c.
	\end{equation}
\end{df}

This definition is actually a very strong definition which requires the sequence to be strictly non-zero. A natural question to ask is why we say that there exists some rational $c > 0$ in which $| a_n | \geq c$ instead of just saying $| a_n | > 0$? Let's look at some examples.

\begin{exmp}{Sequences bounded away from zero vs. Non-zero sequences}{}
	Consider the sequence $\sequence{a_n}{n = m}{\infty}$:
	\begin{equation}
		1, \frac{1}{2}, \frac{1}{3}, \frac{1}{4}, \frac{1}{5}, \dots.
	\end{equation}
	Although each term of the sequence is non-zero, i.e., $| a_n | > 0$ for all $n > m \in \mathbb{N}$, it is not bounded away from zero. You can find a term that's $\epsilon$-close to zero for every $\epsilon$. By \cref{df:reals-bounded_away_from_zero}, there doesn't exist any rational $c$ in which all $| a_n | \geq c$ because for all $c$, there is always a term, $a_{\frac{1}{c} + 1}$ that's smaller than $c$. E.g., if $c = \frac{1}{10}$, then $a_{\frac{1}{\frac{1}{10}} + 1} = a_{11} = \frac{1}{11}$ is smaller than $\frac{1}{10}$.
\end{exmp}

An immediate consequence is that if a sequence is bounded away from zero, it is not equivalent to the zero sequence. This is because the sequence isn't $c$-close to zero.

\section{Ordering the reals}

\section{Exponentiation of a real to a rational number}

\section{The supremum and least upper bound}

\section{Boundedness of sequences}

\section{Convergence}

You may need to grab a cup of tea and think of what this definition actually means. It's important for you to understand it fully. Let it sink in for a while. Well, I'd tell you that this definition is quite a humongous one. So, it'd be better if I break it down into two more digestible chunks: $\epsilon$-boundedness and eventual $\epsilon$-boundedness, in order to ease the understanding.\footnote{These two definitions are not standard in the literature.}

\begin{df}{$(L, \epsilon)$-boundedness}{reals-epsilon_boundedness}
	A sequence of rationals $\sequence{a_n}{n = i}{f}$ is $(L, \epsilon)$-bounded if for all terms in the sequence,
	\begin{equation*}
		|a_n - L| \leq \epsilon.
	\end{equation*}
\end{df}

Expanding the absolute value yields
\begin{gather}
	- \epsilon \leq a_n - L \leq \epsilon \\
	L - \epsilon \leq a_n \leq L + \epsilon,
\end{gather}
i.e., all $a_n$ lies inside $L - \epsilon$ and $L + \epsilon$. Let's consider a sequence of numbers $\sequence{a_n}{n = 0}{\infty}$ illustrated.

\begin{exmp}{$(L, \epsilon)$-bounded sequences}{}
	\begin{enumerate}
		\item The sequence $-1, 1, -1, 1, -1, \dots$ is $(0, 1)$-bounded, but not $(1, 1)$-bounded because $-1$ doesn't lie between $1 - 1 = 0$ and $1 + 1 = 2$.
		\item The sequence $0.9, 0.99, 0.999, 0.9999, 0.99999, \dots$ is $(0.2, 1)$-bounded, but not $(0.001, 1)$-bounded because $0.9$ and $0.99$ doesn't lie between $1 + 0.01 = 1.01$ and $1 - 0.01 = 0.99$.
		\item The sequence $1, \frac{1}{2}, \frac{1}{3}, \frac{1}{4}, \frac{1}{5}, \dots$ is $(0.5, 0.5)$-bounded, but not $(0.5, 0.25)$-bounded because every term that isn't $\frac{1}{2}$, $\frac{1}{3}$ and $\frac{1}{4}$, doesn't lie between $0.5 + 0.25 = 0.75$ and $0.5 - 0.25 = 0.25$.
	\end{enumerate}
\end{exmp}

\begin{df}{Eventual $(L, \epsilon)$-boundedness}{reals-eventual_epsilon_boundedness}
	A sequence of rationals $\sequence{a_n}{n = i}{f}$ is eventually $(L, \epsilon)$-bounded if there exists a number $m$ that lies between $i$ and $f$, such that $\sequence{a_n}{n = m}{f}$ is $(L, \epsilon)$-bounded.
\end{df}

\paragraph{The explanation:} Consider any arbitrary sequence of rationals $\sequence{a_n}{n = 0}{\infty}$ where the limit of $a_n$ is $L$, shown in \cref{fig:reals-convergence_epsilon}. From the definition, there exists some $N$ in which for all $n$ that follows, $|a_n - L| \leq \epsilon$. Expanding the absolute value yields
\begin{gather}
	- \epsilon    \leq a_n - L  \leq + \epsilon   \\
	L - \epsilon  \leq a_n      \leq L + \epsilon.
\end{gather}
I.e., for all $n \geq N$, $a_n$ is between $L - \epsilon$ and $L + \epsilon$. Illustrated in \cref{fig:reals-convergence_epsilon}, it means that there exists some point $a_N$ in which all the terms that follows it lies in the gray rectangle.

\begin{figure}[t]
	\centering
	\includegraphics{realanalysis/realsandsequences/convergence.pdf}
	\caption{Convergence of sequences}
	\label{fig:reals-convergence_epsilon}
\end{figure}
